\documentclass[tikz,border=6pt]{standalone}
\usepackage[UTF8]{ctex}
\usepackage{amsmath}
\usetikzlibrary{arrows.meta,positioning,fit,backgrounds,calc}

\definecolor{cblue}{RGB}{234,243,255}
\definecolor{corange}{RGB}{255,243,232}
\definecolor{cgreen}{RGB}{236,247,238}
\definecolor{cgray}{RGB}{245,246,248}
\definecolor{cline}{RGB}{78,89,104}
\definecolor{ctext}{RGB}{35,42,52}

\tikzset{
	base/.style={draw=cline, rounded corners=2mm, thick, align=center, text=ctext},
	io/.style={base, fill=cgreen, minimum width=23mm, minimum height=8mm, inner sep=3pt},
	block/.style={base, fill=cgray, minimum width=24mm, minimum height=8mm, inner sep=3pt},
	dyn/.style={base, fill=corange, minimum width=31mm, minimum height=11mm, inner sep=4pt},
	kin/.style={base, fill=cblue, minimum width=28mm, minimum height=10mm, inner sep=4pt},
	groupbox/.style={draw=cline!85, dashed, rounded corners=2.5mm, thick, inner sep=5pt},
	arrow/.style={-{Latex[length=2.6mm,width=1.8mm]}, thick, draw=cline},
	fb/.style={-{Latex[length=2.4mm,width=1.7mm]}, thick, draw=cline!85},
	lab/.style={font=\bfseries\scriptsize, text=ctext, fill=white, inner sep=1.5pt}
}

\begin{document}
	\begin{tikzpicture}[node distance=4mm and 5mm, font=\scriptsize]
		
		\node[io] (cmd) {控制输入\\$u_c$};
		\node[block, below=of cmd] (map) {执行器映射\\$u_c \rightarrow \tau$};
		\node[dyn, below=of map] (dyn) {动力学模型\\$M\dot{\nu}+C(\nu)\nu+D(\nu)\nu=\tau+d$};
		\node[io, below=of dyn] (vel) {速度状态\\$\nu=[u,v,r]^T$};
		\node[kin, below=of vel] (kin) {运动学模型\\$\dot{\eta}=J(\psi)\nu$};
		\node[io, below=of kin] (pose) {位姿状态\\$\eta=[x,y,\psi]^T$};
		
		\draw[arrow] (cmd) -- (map);
		\draw[arrow] (map) -- node[right=1mm] {$\tau$} (dyn);
		\draw[arrow] (dyn) -- node[right=1mm] {积分} (vel);
		\draw[arrow] (vel) -- node[right=1mm] {变换} (kin);
		\draw[arrow] (kin) -- node[right=1mm] {积分} (pose);
		
		\node[block, right=8mm of dyn, minimum width=18mm] (dist) {外扰\\$d$};
		\draw[arrow] (dist.west) -- (dyn.east);
		
		\node[block, right=8mm of kin, minimum width=24mm] (ctrl) {导引/控制器\\误差计算与控制律};
%		\draw[fb] (pose.east) -- ++(4mm,0) |- (ctrl.south);
		\draw[fb] (pose.east) -- ++(4mm,0) -| (ctrl.south);
		\draw[fb] (vel.east) -| (ctrl.west);
		\draw[fb] (ctrl.north) |- ++(0,27mm) -| (cmd.east);
		
		\node[align=center, font=\scriptsize, text=ctext] at ($(ctrl)!0.52!(pose.east)+(0.8mm,-2mm)$) {反馈};
		\node[align=center, font=\scriptsize, text=ctext] at ($(ctrl)!0.58!(cmd.east)+(2mm,2mm)$) {期望指令};
		
		\begin{pgfonlayer}{background}
			\node[groupbox, fit=(dyn)(vel), label={[lab]north:动力学层}] {};
			\node[groupbox, fit=(kin)(pose), label={[lab]north:运动学层}] {};
		\end{pgfonlayer}
		
		\node[align=left, text width=7.7cm, anchor=north west, font=\scriptsize, text=ctext]
		at ([yshift=-3mm]pose.south west)
		{\textbf{作用链：}控制输入 $\rightarrow$ 广义力/力矩 $\rightarrow$ 速度变化（动力学） $\rightarrow$ 位姿变化（运动学）。\\
			\textbf{关系概括：}动力学回答“力如何改变速度”，运动学回答“速度如何改变位置与姿态”。};
		
	\end{tikzpicture}
\end{document}